\section{Native Dataset Groupings}
\subsection{OCEAN\_TEMPERATURE\_SALINITY}
\subsubsection{Overview}
\newp
This dataset provides 3D fields of ocean potential temperature and salinity on the lat-lon-cap 90 (llc90) native model grid from the ECCO Version 4 Release 4 (V4r4) ocean and sea-ice state estimate. The dataset is provided on daily-average and monthly-average time resolution. Snapshot data of ocean potential temperature and salinity are also provided.
\\\\
\begin{longtable}{|m{0.15\textwidth}|m{0.64\textwidth}|m{0.12\textwidth}|}
\caption{Coordinates and variables from OCEAN\_TEMPERATURE\_SALINITY (native)}
\label{tab:table-<xarray.Dataset> Size: 47MB
Dimensions:    (time: 1, k: 50, tile: 13, j: 90, i: 90, i_g: 90, j_g: 90,
                k_u: 50, k_l: 50, k_p1: 51, nv: 2, nb: 4)
Coordinates: (12/22)
  * time       (time) datetime64[ns] 8B 2017-12-29T12:00:00
  * k          (k) int32 200B 0 1 2 3 4 5 6 7 8 9 ... 41 42 43 44 45 46 47 48 49
    Z          (k) float32 200B ...
  * tile       (tile) int32 52B 0 1 2 3 4 5 6 7 8 9 10 11 12
  * j          (j) int32 360B 0 1 2 3 4 5 6 7 8 9 ... 81 82 83 84 85 86 87 88 89
  * i          (i) int32 360B 0 1 2 3 4 5 6 7 8 9 ... 81 82 83 84 85 86 87 88 89
    ...         ...
  * k_p1       (k_p1) int32 204B 0 1 2 3 4 5 6 7 8 ... 43 44 45 46 47 48 49 50
    Zp1        (k_p1) float32 204B ...
    time_bnds  (time, nv) datetime64[ns] 16B ...
    Z_bnds     (k, nv) float32 400B ...
    XC_bnds    (tile, j, i, nb) float32 2MB ...
    YC_bnds    (tile, j, i, nb) float32 2MB ...
Dimensions without coordinates: nv, nb
Data variables:
    THETA      (time, k, tile, j, i) float32 21MB ...
    SALT       (time, k, tile, j, i) float32 21MB ...
Attributes: (12/62)
    acknowledgement:                 This research was carried out by the Jet...
    author:                          Ian Fenty and Ou Wang
    cdm_data_type:                   Grid
    comment:                         Fields provided on the curvilinear lat-l...
    Conventions:                     CF-1.8, ACDD-1.3
    coordinates_comment:             Note: the global 'coordinates' attribute...
    ...                              ...
    time_coverage_duration:          P1D
    time_coverage_end:               2017-12-30T00:00:00
    time_coverage_resolution:        P1D
    time_coverage_start:             2017-12-29T00:00:00
    title:                           ECCO Ocean Temperature and Salinity - Da...
    uuid:                            7f3dcd4c-4169-11eb-8b81-0cc47a3f3f51-fields} \\ 
\hline \endfirsthead \endhead \endfoot \hline \endlastfoot
\rowcolor{lightgray} \multicolumn{1}{|c|}{\textbf{Coordinates}} & \multicolumn{1}{|c|}{\textbf{Description of data coordinates}} & \multicolumn{1}{|c|}{\textbf{Unit}}\\ \hline
i &Grid index in x for variables at tracer and 'v' locations &--none--  \\ \hline
i\_g &Grid index in x for variables at 'u' and 'g' locations &--none--  \\ \hline
j &Grid index in y for variables at tracer and 'u' locations &--none--  \\ \hline
j\_g &Grid index in y for variables at 'v' and 'g' locations &--none--  \\ \hline
k &Grid index in z for tracer variables &--none--  \\ \hline
k\_u &Grid index in z corresponding to the bottom face of tracer grid cells ('w' locations) &--none--  \\ \hline
k\_l &Grid index in z corresponding to the top face of tracer grid cells ('w' locations) &--none--  \\ \hline
k\_p1 &Grid index in z for variables at 'w' locations &--none--  \\ \hline
tile &Lat-lon-cap tile index &--none--  \\ \hline
time &Center time of averaging period &--none--  \\ \hline
XC &Longitude of tracer grid cell center &degrees\_east  \\ \hline
YC &Latitude of tracer grid cell center &degrees\_north  \\ \hline
XG &Longitude of 'southwest' corner of tracer grid cell &degrees\_east  \\ \hline
YG &Latitude of 'southwest' corner of tracer grid cell &degrees\_north  \\ \hline
Z &Depth of tracer grid cell center &m  \\ \hline
Zp1 &Depth of tracer grid cell interface &m  \\ \hline
Zu &Depth of the bottom face of tracer grid cells &m  \\ \hline
Zl &Depth of the top face of tracer grid cells &m  \\ \hline
time\_bnds &Time bounds of averaging period &--none--  \\ \hline
XC\_bnds &Longitudes of tracer grid cell corners &--none--  \\ \hline
YC\_bnds &Latitudes of tracer grid cell corners &--none--  \\ \hline
Z\_bnds &Depths of tracer grid cell upper and lower interfaces &--none--  \\ \hline
\rowcolor{lightgray} \multicolumn{1}{|c|}{\textbf{Variables}} & \multicolumn{1}{|c|}{\textbf{Description of data variables}} & \multicolumn{1}{|c|}{\textbf{Unit}}\\ \hline
THETA &Potential temperature  &degree\_C  \\ \hline
SALT &Salinity &1e-3  \\ \hline
\end{longtable}

\newp
\pagebreak
\subsubsection{SALT (native)}
\begin{longtable}{|m{0.06\textwidth}|m{0.3\textwidth}|m{0.45\textwidth}|m{0.12\textwidth}|}
\caption{SALT from OCEAN\_TEMPERATURE\_SALINITY (native)}\label{tab:table-OCEAN_TEMPERATURE_SALINITY_SALT} \\ 
\hline \endhead \hline \endfoot
\rowcolor{lightgray} \textbf{Storage Type} & \textbf{Variable Name} & \textbf{Description} & \textbf{Unit} \\ \hline
float32 & SALT & Salinity & 1e-3 \\ \hline
\multicolumn{4}{|c|}{\cellcolor{lightgray}{\textbf{Common Data format Language (CDL) description}}} \\ \hline
\multicolumn{4}{|c|}{\fontfamily{lmtt}\selectfont{\makecell{\parbox{.95\textwidth}{\vspace*{0.25cm} \footnotesize{float32 SALT(time, k, tile, j, i)\\\
\hspace*{0.5cm}SALT: \_FillValue = 9.96921e+36\\\
\hspace*{0.5cm}SALT: coordinates = YC Z XC time\\\
\hspace*{0.5cm}SALT: coverage\_content\_type = modelResult\\\
\hspace*{0.5cm}SALT: long\_name = Salinity\\\
\hspace*{0.5cm}SALT: standard\_name = sea water salinity\\\
\hspace*{0.5cm}SALT: units = 1e-3\\\
\hspace*{0.5cm}SALT: valid\_max = 41.321231842041016\\\
\hspace*{0.5cm}SALT: valid\_min = 16.73577880859375\\\
}}}}} \\ \hline
\rowcolor{lightgray} \multicolumn{4}{|c|}{\textbf{Comments}} \\ \hline
\multicolumn{4}{|p{1\textwidth}|}{\footnotesize{{Defined using cf convention 'sea water salinity is the salt content of sea water, often on the practical salinity scale of 1978. however, the unqualified term 'salinity' is generic and does not necessarily imply any particular method of calculation. the units of salinity are dimensionless and the units attribute should normally be given as 1e-3 or 0.001 i.e. parts per thousand.' see https://cfconventions.org/data/cf-standard-names/73/build/cf-standard-name-table.html}}} \\ \hline
\end{longtable}

\begin{figure}[H]
\centering
\includegraphics[scale=0.55]{/Users/brucel/ecco/yip/ECCO-Dataset-Production/document_generator/files_general/resource_files/version_specific/v4r4/output_and_granules/output/figures/figures_variables_native/Ocean_Temperature_and_Salinity/SALT.png}
\caption{SALT (native)}
\label{tab:table-OCEAN_TEMPERATURE_SALINITY_SALT-Plot}
\end{figure}
\newpage
\pagebreak
\subsubsection{THETA (native)}
\begin{longtable}{|m{0.06\textwidth}|m{0.3\textwidth}|m{0.45\textwidth}|m{0.12\textwidth}|}
\caption{THETA from OCEAN\_TEMPERATURE\_SALINITY (native)}\label{tab:table-OCEAN_TEMPERATURE_SALINITY_THETA} \\ 
\hline \endhead \hline \endfoot
\rowcolor{lightgray} \textbf{Storage Type} & \textbf{Variable Name} & \textbf{Description} & \textbf{Unit} \\ \hline
float32 & THETA & Potential temperature  & degree\_C \\ \hline
\multicolumn{4}{|c|}{\cellcolor{lightgray}{\textbf{Common Data format Language (CDL) description}}} \\ \hline
\multicolumn{4}{|c|}{\fontfamily{lmtt}\selectfont{\makecell{\parbox{.95\textwidth}{\vspace*{0.25cm} \footnotesize{float32 THETA(time, k, tile, j, i)\\\
\hspace*{0.5cm}THETA: \_FillValue = 9.96921e+36\\\
\hspace*{0.5cm}THETA: coordinates = YC Z XC time\\\
\hspace*{0.5cm}THETA: coverage\_content\_type = modelResult\\\
\hspace*{0.5cm}THETA: long\_name = Potential temperature \\\
\hspace*{0.5cm}THETA: standard\_name = sea water potential temperature\\\
\hspace*{0.5cm}THETA: units = degree C\\\
\hspace*{0.5cm}THETA: valid\_max = 36.425140380859375\\\
\hspace*{0.5cm}THETA: valid\_min = -2.9179372787475586\\\
}}}}} \\ \hline
\rowcolor{lightgray} \multicolumn{4}{|c|}{\textbf{Comments}} \\ \hline
\multicolumn{4}{|p{1\textwidth}|}{\footnotesize{{Sea water potential temperature is the temperature a parcel of sea water would have if moved adiabatically to sea level pressure. note: the equation of state is a modified unesco formula by jackett and mcdougall (1995), which uses the model variable potential temperature as input assuming a horizontally and temporally constant pressure of \$p\_0=-g 
ho\_\{0\} z\$.}}} \\ \hline
\end{longtable}

\begin{figure}[H]
\centering
\includegraphics[scale=0.55]{/Users/brucel/ecco/yip/ECCO-Dataset-Production/document_generator/files_general/resource_files/version_specific/v4r4/output_and_granules/output/figures/figures_variables_native/Ocean_Temperature_and_Salinity/THETA.png}
\caption{THETA (native)}
\label{tab:table-OCEAN_TEMPERATURE_SALINITY_THETA-Plot}
\end{figure}
\newpage